\section{Method}
\label{sec:method}

\subsection{Background: The Monitored Detector}
\label{sec:background-detector}

The detector this paper monitors is a camera/LiDAR BEV fusion architecture; we summarize only the components this paper's monitoring layer directly depends on, since the monitor itself is detector-agnostic and this paper's own contribution begins with the monitoring layer described in the rest of this section.

Camera and LiDAR streams are each encoded into a shared bird's-eye-view (BEV) feature grid, $F_{\mathrm{cam}}, F_{\mathrm{lid}} \in \mathbb{R}^{H_b \times W_b \times C}$. A Cross-modal Consistency Probe (CCP) computes, per BEV grid cell $(i,j)$, a consistency score
\begin{equation}
S(i,j) = \sigma\Big(\mathrm{MLP}_S\big[F_{\mathrm{cam}}(i,j) \oplus F_{\mathrm{lid}}(i,j) \oplus |F_{\mathrm{cam}}(i,j) - F_{\mathrm{lid}}(i,j)|\big]\Big) \in [0,1],
\label{eq:ccp}
\end{equation}
from the concatenated ($\oplus$) camera and LiDAR feature vectors at that cell and their absolute difference, so that $1 - S(i,j)$ is a cross-modal disagreement signal intended to flag localized geometric or semantic mismatch between the two modalities. A Gated Adaptive Fusion Module (GAFM) uses $S$ to weight each modality's contribution to the fused BEV representation at the grid-cell level, and an Adversarial Consistency Training (ACT) objective trains the detector, the CCP, and the fusion gate jointly under a loss combining detection accuracy, a CCP contrastive-consistency term, and an attribution-regularization term, with adversarially perturbed camera and LiDAR branches included in the training loop to encourage $S$ to remain informative under perturbation. This paper's own contribution begins with the next subsection: everything from here on -- the conformal calibration, the sequential test, the spatial multiplicity control, and the use (and, as Section~\ref{sec:experiments} reports, the eventual disuse) of $S$ as a monitoring covariate -- is this paper's method, evaluated against real data in Section~\ref{sec:experiments}.

\begin{figure}[t]
\centering
\begin{tikzpicture}[
    box/.style={draw, rounded corners, align=center, minimum height=8mm, text width=1.55cm, inner sep=2pt, font=\scriptsize},
    wbox/.style={draw, rounded corners, align=center, minimum height=7mm, text width=3.55cm, inner sep=2pt, font=\scriptsize},
    arr/.style={-{Stealth[length=1.6mm]}, thick}
  ]
  \node[box] (frame) at (0, 0) {Frame $t$\\(image, LiDAR)};
  \node[box] (det) at (2.1, 0) {Detector\\$f_\theta$ (background)};
  \node[box] (conf) at (0, -1.5) {Conformal\\$m(t)$ vs.\ $\hat{q}_\alpha$};
  \node[box] (ccp) at (2.1, -1.5) {CCP $S_t$\\disagree.\ $\overline{1-S_t}$};
  \node[wbox] (bet) at (1.05, -3.0) {Betting rule $\lambda_t$\\(blind or CCP-informed)};
  \node[box] (wealth) at (0, -4.5) {Wealth\\$K_t$};
  \node[box] (alarm) at (2.1, -4.5) {Alarm if\\$K_t \ge 1/\delta$};

  \draw[arr] (frame) -- (det);
  \draw[arr] (det) -- (conf);
  \draw[arr] (det) -- (ccp);
  \draw[arr] (conf) -- (bet) node[midway, left, font=\tiny] {$m(t)$};
  \draw[arr] (ccp) -- (bet) node[midway, right, font=\tiny] {$\overline{1-S_t}$};
  \draw[arr] (bet) -- (wealth);
  \draw[arr] (wealth) -- (alarm);
\end{tikzpicture}
\caption{Per-frame monitoring pipeline. The detector $f_\theta$ (dashed conceptually from the rest, since it is background infrastructure -- Section~\ref{sec:background-detector}) produces box predictions and the CCP score $S_t$; everything downstream of it is this paper's own contribution. The spatial variant (Section~\ref{sec:spatial}) replicates the conformal-through-wealth path once per BEV cell-cluster. The betting rule is either covariate-blind (uses only the $m(t)$ path) or CCP-informed (also uses the disagreement path).}
\label{fig:pipeline}
\end{figure}

Figure~\ref{fig:pipeline} sketches the full per-frame pipeline before we define each stage precisely.

\subsection{Problem Setup}
\label{sec:setup}

We observe a stream of driving frames indexed by $t = 1, 2, \ldots$, one per LiDAR sweep. At every $t$, the base 3D detector $f_\theta$ described above produces per-object box predictions from the frame's fused camera/LiDAR BEV features; the monitor itself is detector-agnostic and does not depend on $f_\theta$'s specific architecture. We wrap those predictions in split conformal prediction, calibrated once on a held-out set of frames drawn from the nominal (least-degraded) regime available in the data. For object $i$ in frame $t$, the nonconformity score is the $\ell_1$ box-regression residual restricted to grid cells a ground-truth object actually occupies,
\[
s_i(t) = \big\| \hat{B}_i(t) - B_i(t) \big\|_1 ,
\]
and the calibrated threshold $\hat{q}_\alpha$ is the usual finite-sample-valid split-conformal quantile: the $\lceil (n+1)(1-\alpha) \rceil / n$ empirical quantile of the calibration set's nonconformity scores $\{s_i\}_{i=1}^{n}$, so that a new object is covered iff $s_i(t) \le \hat{q}_\alpha$. The frame-level miscoverage rate is
\[
m(t) = \frac{1}{n_t} \sum_{i=1}^{n_t} \mathbf{1}\{s_i(t) > \hat{q}_\alpha\},
\]
with $m(t) := 0$ defined for frames with no matched objects, since an empty frame carries no evidence of miscoverage. The null hypothesis the monitor tests is that the pipeline is still operating nominally, $H_0: \mathbb{E}[m(t) \mid \mathcal{F}_{t-1}] \le \alpha$ for all $t$; the goal is to detect the first violation -- weather-onset degradation -- with an anytime-valid false-alarm guarantee, not merely a guarantee at one fixed check time.

\subsection{The e-Process and Ville's Inequality}
\label{sec:eprocess}

The wealth process is
\begin{align}
K_0 &= 1, \nonumber \\
K_t &= \prod_{s=1}^{t} \big(1 + \lambda_s (m(s) - \alpha)\big), \qquad \lambda_s \in \left[0, \tfrac{1}{\alpha}\right] \text{ chosen } \mathcal{F}_{s-1}\text{-measurably.}
\label{eq:wealth}
\end{align}
We work through why this construction gives an anytime-valid test, rather than only stating the result, since the two steps -- the supermartingale property and Ville's inequality -- are each short but easy to state imprecisely.

\paragraph{Step 1: $K_t$ is a nonnegative supermartingale under $H_0$.}
Non-negativity requires $1 + \lambda_t(m(t) - \alpha) \ge 0$ for every realizable $m(t) \in [0,1]$. Since $\lambda_t \ge 0$, the expression is affine and non-decreasing in $m(t)$, so its minimum over $m(t) \in [0,1]$ occurs at $m(t) = 0$, giving $1 - \lambda_t \alpha$. This is non-negative exactly when $\lambda_t \le 1/\alpha$ -- the bound our implementation enforces on every betting rule (\texttt{WealthProcess.lambda\_max} in \texttt{conformal\_monitor/betting.py}), and the reason the bound is $1/\alpha$ specifically, not the also-plausible-looking $1/(1-\alpha)$ (which only bounds the opposite endpoint, $m(t)=1$, and is too loose at the endpoint that actually binds). For the supermartingale property itself, $\lambda_t$ is chosen $\mathcal{F}_{t-1}$-measurably -- it may depend on the entire history up to and including frame $t-1$, but not on $m(t)$ itself -- so, conditioning on $\mathcal{F}_{t-1}$,
\[
\mathbb{E}\big[K_t \mid \mathcal{F}_{t-1}\big] = K_{t-1} \cdot \mathbb{E}\big[1 + \lambda_t(m(t) - \alpha) \mid \mathcal{F}_{t-1}\big] = K_{t-1}\big(1 + \lambda_t(\mathbb{E}[m(t) \mid \mathcal{F}_{t-1}] - \alpha)\big),
\]
using that $\lambda_t$ and $K_{t-1}$ are both $\mathcal{F}_{t-1}$-measurable and can be pulled outside the conditional expectation. Under $H_0$, $\mathbb{E}[m(t) \mid \mathcal{F}_{t-1}] \le \alpha$, so the bracketed term is $\le 1$ whenever $\lambda_t \ge 0$ (which it always is, by construction), giving $\mathbb{E}[K_t \mid \mathcal{F}_{t-1}] \le K_{t-1}$: precisely the supermartingale inequality, holding for \emph{every} choice of predictable, non-negative $\lambda_t \le 1/\alpha$, which is exactly the freedom the different betting rules in Sections~\ref{sec:blind}--\ref{sec:informed} exploit without ever needing to reprove this step.

\paragraph{Step 2: Ville's inequality bounds the running maximum.}
For a nonnegative supermartingale $K_t$ with $K_0 = 1$, Ville's maximal inequality gives, for any $c > 0$,
\[
\mathbb{P}_{H_0}\!\left(\sup_{t \ge 0} K_t \ge c\right) \le \frac{\mathbb{E}[K_0]}{c} = \frac{1}{c}.
\]
Setting $c = 1/\delta$ gives $\mathbb{P}_{H_0}\!\left(\sup_t K_t \ge 1/\delta\right) \le \delta$. The alarm rule -- reject $H_0$ the first time $K_t \ge 1/\delta$ -- therefore has false-alarm probability at most $\delta$ \emph{regardless of the (random) stopping time at which the process is examined}: the guarantee is over the entire trajectory's supremum, not over $K_t$ at one pre-specified $t$, which is exactly what makes checking continuously, with no fixed monitoring horizon and no correction for how many times the test has been evaluated, statistically valid rather than an uncorrected multiple-testing error. This two-step argument -- the supermartingale property from Step 1, Ville's inequality from Step 2 -- is the established backbone of testing by betting and is not this paper's contribution; what varies across the rest of this section, and what this paper's own contribution actually depends on, is how the predictable sequence $\lambda_t$ is chosen.

\subsubsection{A Worked Numeric Example}
\label{sec:worked-example}
To build intuition before the general betting rules, we trace $K_t$ over five synthetic frames by hand, with $\alpha = 0.2$ (so $\lambda_{\max} = 5$) and a fixed, non-adaptive $\lambda_t = 2$ for simplicity (an actual bettor updates $\lambda_t$ per Sections~\ref{sec:blind}--\ref{sec:informed}; holding it fixed here isolates what the recursion itself does). Suppose the observed miscoverage sequence is $m = (0.1, 0.15, 0.6, 0.55, 0.5)$ -- nominal for the first two frames, then abruptly degraded from frame 3 onward, mimicking a weather onset at $t=3$:
\begin{align*}
K_1 &= K_0 \big(1 + 2(0.10 - 0.2)\big) = 1 \times 0.80 = 0.80, \\
K_2 &= K_1 \big(1 + 2(0.15 - 0.2)\big) = 0.80 \times 0.90 = 0.72, \\
K_3 &= K_2 \big(1 + 2(0.60 - 0.2)\big) = 0.72 \times 1.80 = 1.296, \\
K_4 &= K_3 \big(1 + 2(0.55 - 0.2)\big) = 1.296 \times 1.70 = 2.203, \\
K_5 &= K_4 \big(1 + 2(0.50 - 0.2)\big) = 2.203 \times 1.60 = 3.525.
\end{align*}
Two things are worth noticing in this trace. First, wealth \emph{shrinks} while $m(t) < \alpha$ (frames 1--2): a bettor staking $\lambda_t > 0$ on ``miscoverage will exceed $\alpha$'' loses money, by construction, when the pipeline is behaving nominally -- this is the price of testing at all, and is what keeps the process a supermartingale rather than an unconditionally growing one. Second, wealth grows multiplicatively, not additively, once $m(t) > \alpha$ consistently (frames 3--5): each factor $> 1$ compounds against the previous wealth, which is why a real onset that persists for several frames can cross a large threshold $1/\delta$ (e.g.\ $1/\delta = 20$ for $\delta = 0.05$) well before $t \to \infty$, even though no single frame's factor is large by itself. In this toy trace, $K_5 = 3.525$ would already trigger an alarm at $\delta = 0.3$ ($1/\delta \approx 3.33$) but not yet at $\delta = 0.1$ or $\delta = 0.05$ -- illustrating concretely why Table~\ref{tab:snowy-results}'s detection delay lengthens as $\delta$ shrinks: a smaller false-alarm budget requires more accumulated evidence (a higher wealth threshold) before the anytime-valid guarantee permits an alarm.

\subsection{Covariate-Blind Baselines}
\label{sec:blind}

Two covariate-blind betting rules choose $\lambda_t$ purely from the history of $m(t)$ itself, reproducing~\cite{MonroyMunoz2026WACV} as the baseline to beat. Approximate GRAPA, with $X_t = m(t) - \alpha$, sets
\begin{align*}
\lambda_t = \operatorname{clip}\!\left(\frac{\overline{X}_{<t}}{\mathbb{E}[X_{<t}^2] + \varepsilon}, \, 0, \, \lambda_{\max}\right),
\end{align*}
tracking the running first and second moments of the centered miscoverage. Scale-Free Online Gradient Descent instead performs AdaGrad-style updates on the log-wealth gradient, with $g_t = X_t / (1 + \lambda_t X_t)$ and $\eta_t = \lambda_{\max} / \sqrt{1 + \sum_{s \le t} g_s^2}$,
\begin{align*}
\lambda_{t+1} = \operatorname{clip}\!\left(\lambda_t + \eta_t \, g_t, \, 0, \, \lambda_{\max}\right).
\end{align*}
Both are exact reimplementations of the adaptive betting rules \cite{MonroyMunoz2026WACV} uses, run here against a different sensing modality and detection task than the 2D visual tracking they were originally evaluated on.

\subsection{Covariate-Informed Betting}
\label{sec:informed}

Alongside the two contributions in Sections~\ref{sec:spatial} and~\ref{sec:experiments}, we also investigated conditioning the betting rate on an external signal available at the same timestep as $m(t)$ but computed from independent evidence: the disagreement output of the monitored detector's own CCP, $1-S$ (Equation~\ref{eq:ccp}). We describe the mechanism in full because Section~\ref{sec:experiments} reports, in equal detail, why it did not deliver the advantage motivating it -- an outcome worth understanding precisely rather than skipping past. Our covariate-informed bettor wraps either covariate-blind base bettor and scales its stake by the frame's mean disagreement,
\[
\lambda_t = \operatorname{clip}\!\Big(\lambda_t^{\text{base}} \cdot \big(1 + \kappa \cdot \overline{(1 - S_t)}\big), \, 0, \, \lambda_{\max}\Big),
\]
where $\kappa \ge 0$ gates the covariate's influence ($\kappa = 0$ recovers the base bettor exactly) and $\overline{(1-S_t)}$ is the frame- (or cell-)averaged disagreement. The intended mechanism is that a scene beginning to degrade shows rising cross-modal disagreement before enough individual object detections have actually missed coverage for $m(t)$ to move -- disagreement as a leading indicator, miscoverage as the lagging one the covariate-blind rules are stuck reacting to. Section~\ref{sec:experiments} reports what we found when we tested that mechanism against real data: two distinct failures, neither in the betting-rule mathematics itself, that between them mean this specific mechanism does not yet deliver the advantage it is designed to.

\subsection{Spatially Resolved Monitoring with Online Multiplicity Control}
\label{sec:spatial}

A single global wealth process answers only "has anything, anywhere, gone wrong" -- not useful to a downstream planner that needs to know \emph{which} region of the scene to distrust. We instead partition the BEV grid into an $(n_h \times n_w)$ layout of cell-clusters (average-pooling per-cell miscoverage and disagreement down to this resolution) and run one independent sequential tester per cluster, each with its own wealth process and its own instance of whichever bettor is under test. Testing $n = n_h n_w$ hypotheses simultaneously, every frame, requires multiplicity control or the overall false-alarm rate inflates with $n$; we support two regimes. Bonferroni allocates each cell a fixed budget $\delta/n$ and, once a cell alarms, marks it permanently flagged -- valid under arbitrary dependence but conservative, and the sticky-once-alarmed semantics is what we use for the per-cell detection-delay evaluation regardless of which regime produced the alarm. The e-BH alternative treats each cell's current wealth $K_t^{(i)}$ as an e-value: for a nonnegative supermartingale started at $K_0=1$, $\mathbb{E}_{H_0}[K_t^{(i)}] \le \mathbb{E}_{H_0}[K_0^{(i)}] = 1$ at any fixed $t$ (immediate from Section~\ref{sec:eprocess}'s Step 1, since the supermartingale inequality applied recursively from $t=0$ gives $\mathbb{E}[K_t] \le \mathbb{E}[K_{t-1}] \le \cdots \le K_0 = 1$) -- exactly the definition of an e-value under $H_0^{(i)}$: a nonnegative random variable with expectation at most 1. Given $n = n_h n_w$ such e-values at a fixed $t$, one per cell, e-BH~\cite{WangRamdas2022JRSSB} rejects the set
\[
R_t = \Big\{ i : K_t^{(i)} \ge \frac{n}{k^\ast q} \Big\}, \qquad k^\ast = \max\Big\{ k \in \{0, \ldots, n\} : e_{(k)} \ge \frac{n}{k\,q} \Big\},
\]
where $e_{(1)} \ge e_{(2)} \ge \cdots \ge e_{(n)}$ is the descending order statistic of the $n$ current wealths, and $k^\ast = 0$ gives $R_t = \emptyset$. The rule is a direct e-value analogue of Benjamini-Hochberg: it thresholds the $k$-th largest e-value against a level that grows with $k$ (more permissively for larger candidate rejection sets), then commits to the \emph{largest} such $k$ -- the same "largest crossing index" structure as classical BH's $p$-value threshold, but with the inequality direction reversed because e-values are large, not small, under evidence against the null. Its false discovery rate guarantee, $\mathbb{E}[\mathrm{FDR}_t] \le q$ at every fixed $t$, holds under \emph{arbitrary dependence} among the $n$ e-values -- unlike classical BH on $p$-values, which requires independence or specific positive-dependence structures (PRDS) to control FDR without an extra $\log n$ correction factor -- because the e-value definition ($\mathbb{E}[e_i] \le 1$ under $H_0^{(i)}$) is closed under averaging regardless of the joint dependence structure of the $e_i$, which is precisely the guarantee needed here: spatially adjacent BEV cells' miscoverage and disagreement signals are driven by overlapping physical causes (a single snow squall affects many neighboring cells at once) and are certainly not independent. Unlike Bonferroni's sticky once-alarmed flags, e-BH is applied fresh at every timestep from each cell's current wealth, so it produces a live, non-sticky map that can un-flag a cell once its local evidence subsides -- appropriate for a spatial map meant to reflect \emph{current}, not historical, trustworthiness. Algorithm~\ref{alg:spatial-monitor} summarizes the per-frame update.

\subsubsection{Betting-Rule Comparison: Theoretical Properties and Tradeoffs}
\label{sec:betting-comparison}
Table~\ref{tab:betting-comparison} summarizes the three betting rules evaluated in this paper, all valid choices of the predictable $\lambda_t$ in Equation~\ref{eq:wealth} (so all three inherit Section~\ref{sec:eprocess}'s anytime-valid guarantee unconditionally; what differs is statistical power, not validity). AGRAPA is a plug-in estimator of the (asymptotically) growth-rate-optimal constant bet, tracking the running first and second moments of $X_t = m(t)-\alpha$ and converging quickly when the true miscoverage rate is close to stationary, but reacting relatively slowly to an abrupt regime change, since a single new observation only nudges a running average computed over all prior history. SF-OGD instead takes an online-convex-optimization view, performing AdaGrad-style gradient ascent directly on (an approximation to) the log-wealth objective with a decaying step size; this makes it more reactive to recent observations by construction (the step size $\eta_t$ shrinks as $1/\sqrt{t}$, but each step still moves $\lambda_t$ in the direction the \emph{most recent} gradient indicates), at the cost of higher variance in $\lambda_t$ itself, particularly early in a stream before $\sum_{s\le t} g_s^2$ has accumulated enough to stabilize $\eta_t$. Both are exact reimplementations of the rules \cite{MonroyMunoz2026WACV} uses, run here against a different sensing modality and detection task than the 2D visual tracking they were originally evaluated on. The CCP-informed variant wraps either base rule multiplicatively rather than replacing it, so it inherits whichever base rule's reactivity/stability tradeoff was chosen, and adds a single scalar gain $\kappa$ controlling how strongly the covariate perturbs the base bet; $\kappa=0$ recovers the base rule exactly, which is why Section~\ref{sec:experiments} can report the covariate-informed and covariate-blind results as a controlled comparison (same base rule, same wealth-process mechanics, differing only in whether $\overline{1-S_t}$ enters the stake) rather than two unrelated procedures.

\begin{table}[t]
\centering
\caption{Betting rules compared in this paper. All three are valid choices of $\lambda_t$ in Equation~\ref{eq:wealth} and inherit the same anytime-valid guarantee unconditionally (Section~\ref{sec:eprocess}); they differ in statistical power, not validity.}
\label{tab:betting-comparison}
\begin{threeparttable}
\begin{tabular}{@{}p{1.6cm}p{2.6cm}p{2.6cm}@{}}
\toprule
Rule & Mechanism & Reactivity / stability \\
\midrule
AGRAPA & running 1st/2nd moment plug-in & smooth, slower to react to abrupt shift \\
SF-OGD & AdaGrad-style gradient ascent on log-wealth & more reactive, higher early-stream variance \\
CCP-informed & multiplicative covariate gain $\kappa$ on either base rule & inherits base rule's tradeoff; adds covariate sensitivity (Section~\ref{sec:experiments} tests whether this helps) \\
\bottomrule
\end{tabular}
\end{threeparttable}
\end{table}

\begin{algorithm}[t]
\caption{Per-frame spatial monitoring update}
\label{alg:spatial-monitor}
\footnotesize
\begin{algorithmic}[1]
\Require frame $t$; detector output; calibrated quantile $\hat{q}_\alpha$; grid shape $(n_h, n_w)$; correction mode
\State compute per-cell nonconformity, coverage, and miscoverage $m^{(i)}(t)$ for each cluster $i$
\State compute per-cell CCP disagreement $\overline{(1-S)}^{(i)}(t)$
\For{each cluster $i = 1, \ldots, n_h n_w$}
    \State $\lambda_t^{(i)} \gets$ bettor.\Call{next\_lambda}{$\overline{(1-S)}^{(i)}(t)$}
    \State $K_t^{(i)} \gets K_{t-1}^{(i)} \cdot \big(1 + \lambda_t^{(i)}(m^{(i)}(t) - \alpha)\big)$
    \State bettor.\Call{update}{$m^{(i)}(t)$, $\overline{(1-S)}^{(i)}(t)$}
\EndFor
\If{correction = Bonferroni}
    \State flag cell $i$ (permanently) if $K_t^{(i)} \ge n_h n_w / \delta$
\Else \Comment{e-BH}
    \State flag the top-$k^\ast$ cells by $K_t^{(i)}$, $k^\ast$ = largest $k$ with $K_{(k)} \ge (n_h n_w) / (k\,\delta)$
\EndIf
\State \Return flagged-cell map (untrustworthy regions at time $t$)
\end{algorithmic}
\end{algorithm}
